Cryptographic protections such as symmetric encryption (\Cref{sec:symmetric_cryptography}) and digital signatures (\Cref{sec:digital_signatures}) are more likely to be bypassed than cracked. In fact, as Kerckhoffs's principle entails (\Cref{sec:cryptography_basics}), the problem of protecting data and services in the realm of applied cryptography often reduces to the problem of guaranteeing \textbf{the secure management of keys} throughout their lifecycle.

The term \textbf{key management} comprises all activities concerning the handling of keys and related information throughout their lifecycle, including their generation, use, revocation, and destruction.

\readBlock{More information on key management}{The \gls{nist}'s series of special publications on key management \cite{barker_recommendation_2020_p1,barker_recommendation_2019_p2,barker_recommendation_2015} are an excellent --- albeit elaborate --- source of more information on key management. Note that a non-negligible part of the content below is taken from such special publications.}